\subsection*{TEXT / SYSTEM}
\begin{longtable}{@{}p{1.7cm}p{5.2cm}p{7.3cm}@{}}
\toprule
\textbf{Task} & \textbf{Parameters} & \textbf{Action} \\
\midrule
\endfirsthead
\toprule
\textbf{Task} & \textbf{Parameters} & \textbf{Action} \\
\midrule
\endhead
\ttfamily 0 & D1.W = length & Display string at (A1), D1.W bytes, then CR/LF \\
\ttfamily 1 & D1.W = length & Display string at (A1), D1.W bytes, no CR/LF \\
\ttfamily 2 & D1.L (out) = length & Read a string into the buffer at (A1); returns length read \\
\ttfamily 3 & D1.L = value & Display D1.L as a signed decimal number \\
\ttfamily 4 & D1.L (out) & Read a number (prompted) into D1.L \\
\ttfamily 5 & D1.B (out) & Read a single character (prompted) into D1.B \\
\ttfamily 6 & D1.B = char & Display a single character \\
\ttfamily 7 & D1.B (out) & Set D1.B to 1 if keyboard input is pending, else 0 \\
\ttfamily 8 & D1.L (out) & Return time of day: hundredths of a second since midnight \\
\ttfamily 9 & — & Terminate the program (SIMHALT) \\
\ttfamily 10 & A1 = text ptr & Print text at (A1) — unavailable in a browser, logged only \\
\ttfamily 11 & D1.W: \$FF00=clear, \$00FF=get pos (out), else low/high byte = col/row & Position the text cursor, clear the screen, or read the cursor position \\
\ttfamily 12 & D1.B = 0 (off) / non-zero (on) & Keyboard echo on/off \\
\ttfamily 13 & A1 = NUL-terminated string & Display NUL-terminated string at (A1), then CR/LF \\
\ttfamily 14 & A1 = NUL-terminated string & Display NUL-terminated string at (A1), no CR/LF \\
\ttfamily 15 & D1.L = value, D2.B = radix & Display D1 converted to the given radix (e.g. 16 for hex) \\
\ttfamily 16 & D1.B = 0 (off) / non-zero (on) & Input prompt on/off \\
\ttfamily 17 & A1 = NUL-terminated string, D1.L = value & Display NUL-terminated string at (A1), no CR/LF, then display D1.L as a decimal number, no CR/LF \\
\ttfamily 18 & A1 = NUL-terminated string, D1.L (out) & Display NUL-terminated string at (A1), no CR/LF, then read a number (prompted) into D1.L \\
\ttfamily 19 & D1.L = 0 for last key codes, or up to 4 packed key codes to test; D1.L (out) = result & If D1.L in is 0, return the last key-up code (bits 16-23) and key-down code (bits 0-7); otherwise test up to four packed key codes and return \$FF in each byte whose key is currently held down \\
\ttfamily 20 & D1.L = value, D2.B = field width & Display D1.L as a zero/space-padded decimal number \\
\ttfamily 21 & D1.L = BGR color, D2.L packed: bits 0-15 (D2.W) = style bits (bold=1, italic=2, underline=4, strike=8), bits 16-23 = size in points, bits 24-31 = font id (unused) & Set the text font color, style, and size - the bitmap font only supports integer scale steps, so size is rounded to the nearest multiple of the native 8pt size. Load D2 as a full .L - the size lives above bit 15, so a plain MOVE.W into D2 would leave it untouched \\
\ttfamily 22 & D1.L packed (in): bits 16-31 = row, bits 0-15 = col; D1.B (out) = char & Return the character currently at the given row/col \\
\ttfamily 23 & D1.L = hundredths of a second & Delay - a real wall-clock pause while Run is active \\
\ttfamily 24 & D1.B & Enable/disable shortcut keys - no effect (no desktop-style shortcuts to suppress in a web port) \\
\ttfamily 25 & D1.L packed: bits 16-31 = row, bits 0-15 = col; D2.L packed: bits 16-31 = height, bits 0-15 = width; D3.W = direction (0=up, 1=down) & Scroll a text rectangle up or down \\
\ttfamily 30 & — & Clear the instruction cycle counter \\
\ttfamily 31 & D1.L (out) & Return the instruction cycle counter \\
\ttfamily 32.0 & D1.B(in) = 0 & Display the hardware window (7-segment, LEDs, switches) \\
\ttfamily 32.1 & D1.B(in) = 1; D1.L(out) = address & Return the memory address of the 7-segment display \\
\ttfamily 32.2 & D1.B(in) = 2; D1.L(out) = address & Return the memory address of the LEDs \\
\ttfamily 32.3 & D1.B(in) = 3; D1.L(out) = address & Return the memory address of the toggle switches \\
\ttfamily 32.4 & D1.B(in) = 4; D1.L(out) = version & Return the simulator version number \\
\ttfamily 32.5 & D1.B(in) = 5 & Enable exception processing \\
\ttfamily 32.6 & D1.B(in) = 6; D2.B = mask (bit 7 = enable, bits 0-6 = IRQ level), D3.L = period in ms & Set up (or cancel) auto-repeating IRQ generation for a hardware level - same underlying state as the Hardware window's own "Auto every" checkbox/interval for that level \\
\ttfamily 32.7 & D1.B(in) = 7; D1.L(out) = address & Return the memory address of the push-button switches \\
\ttfamily 33 & D1.L=1 windowed, D1.L=2 fullscreen, D1.L=W<<16|H set size, D1.L=0 get size (out, packed the same way) & Get/set the output window mode and size (fullscreen via the real browser Fullscreen API - Esc, or the window's own maximize/restore button, to exit) \\
\bottomrule
\end{longtable}
\subsection*{FILE I/O}
\begin{longtable}{@{}p{1.7cm}p{5.2cm}p{7.3cm}@{}}
\toprule
\textbf{Task} & \textbf{Parameters} & \textbf{Action} \\
\midrule
\endfirsthead
\toprule
\textbf{Task} & \textbf{Parameters} & \textbf{Action} \\
\midrule
\endhead
\ttfamily 50 & D0.W(out) = result & Close all open files \\
\ttfamily 51 & A1 = filename, D1.L(out) = handle, D0.W(out) = result & Open an existing file for read/update (falls back to read-only if it isn't writable) \\
\ttfamily 52 & A1 = filename, D1.L(out) = handle, D0.W(out) = result & Create a new file (or truncate an existing one) for read/write \\
\ttfamily 53 & D1.L=handle, A1=buffer, D2.L(in)=bytes to read, D2.L(out)=bytes actually read, D0.W(out)=result & Read from a file \\
\ttfamily 54 & D1.L=handle, A1=buffer, D2.L=bytes, D0.W(out)=result & Write to a file \\
\ttfamily 55 & D1.L=handle, D2.L=absolute byte offset from the start of the file, D0.W(out)=result & Position (seek) within a file \\
\ttfamily 56 & D1.L = handle, D0.W(out) = result & Close a file \\
\ttfamily 57 & A1 = filename, D0.W(out) = result & Delete a file \\
\ttfamily 58 & D1.L(in)=dialog mode (0=Open, else Save - label text only, behavior is the same either way), D1.L(out)=1 if a file was picked/0 if canceled, A1=title string ptr (or 0), A2=extension filter string ptr (or 0, not used to filter here), A3=path buffer (in: initial value; out: chosen filename) & Show a file picker - implemented as a prompt() listing the files already in the virtual file system, since a browser can't show a native OS file dialog \\
\ttfamily 59 & A1 = filename, D0.W(out) = result & Check whether a file exists and is writable, without opening it (D1 has no effect in this build - the real task's mode/sub-operation switch was never implemented on either the original or this port's file-op routine, which unconditionally just probes the file) \\
\bottomrule
\end{longtable}
\par\noindent\textit{Every file task above returns its result code in D0.W: 0 = success, 1 = end-of-file (task 53 only), 2 = error, 3 = opened read-only (tasks 51/59 only, when the file exists but couldn't be opened for writing).}\par\medskip
\subsection*{MOUSE / KEYBOARD}
\begin{longtable}{@{}p{1.7cm}p{5.2cm}p{7.3cm}@{}}
\toprule
\textbf{Task} & \textbf{Parameters} & \textbf{Action} \\
\midrule
\endfirsthead
\toprule
\textbf{Task} & \textbf{Parameters} & \textbf{Action} \\
\midrule
\endhead
\ttfamily 60 & D1.W packed: bits 8-15 = IRQ level 0-7, bits 0-2 = arm button-down/button-up/move & Arm the mouse button-down/up/move IRQ(s) at the given level (a level >7 is ignored) - click the graphics screen to give it focus first \\
\ttfamily 61 & D1.W(out)=X, D2.W(out)=Y, D3.B(out)=buttons & Read the current mouse position and button state \\
\ttfamily 62 & D1.W packed: bits 8-15 = IRQ level 0-7, bits 0-1 = arm key-down/key-up & Arm the keyboard key-down/up IRQ(s) at the given level (a level >7 is ignored) - click the graphics screen to give it focus first \\
\bottomrule
\end{longtable}
\subsection*{SOUND}
\begin{longtable}{@{}p{1.7cm}p{5.2cm}p{7.3cm}@{}}
\toprule
\textbf{Task} & \textbf{Parameters} & \textbf{Action} \\
\midrule
\endfirsthead
\toprule
\textbf{Task} & \textbf{Parameters} & \textbf{Action} \\
\midrule
\endhead
\ttfamily 70 & A1 = WAV filename, D0.W(out) = result & Play a WAV file directly, one-shot (from the uploaded virtual filesystem, via Web Audio) \\
\ttfamily 71 & A1 = WAV filename, D1.B(in) = slot to load it into & Load a WAV into a player slot you choose, for fast replay later via task 72 (no result code is returned by this build) \\
\ttfamily 72 & D1.B(in) = slot, D0.W(out) = result & Play a previously-loaded WAV from the given slot \\
\ttfamily 73 & A1 = WAV filename, D0.W(out) = result & Play a WAV file directly, one-shot ("DirectX" player variant) \\
\ttfamily 74 & A1 = WAV filename, D1.B(in) = slot to load it into, D0.W(out) = result & Load a WAV into a "DirectX" player slot you choose \\
\ttfamily 75 & D1.B(in) = slot, D0.W(out) = result & Play a previously-loaded WAV from the given "DirectX" slot \\
\ttfamily 76 & D1.B = slot, D2.L = command, D0.W(out) = result & Control the standard player for a loaded slot (stop/pause/volume, depending on the D2 command value) \\
\ttfamily 77 & D1.B = slot, D2.L = command, D0.W(out) = result & Control the "DirectX" player for a loaded slot (stop/pause/volume, depending on the D2 command value) \\
\bottomrule
\end{longtable}
\subsection*{GRAPHICS}
\begin{longtable}{@{}p{1.7cm}p{5.2cm}p{7.3cm}@{}}
\toprule
\textbf{Task} & \textbf{Parameters} & \textbf{Action} \\
\midrule
\endfirsthead
\toprule
\textbf{Task} & \textbf{Parameters} & \textbf{Action} \\
\midrule
\endhead
\ttfamily 80 & D1.L = BGR color & Set the pen (line/outline) color \\
\ttfamily 81 & D1.L = BGR color & Set the fill color \\
\ttfamily 82 & D1.W = X, D2.W = Y & Draw a single pixel \\
\ttfamily 83 & D1.W=X, D2.W=Y, D1.L(out)=color & Read the color of a pixel \\
\ttfamily 84 & D1.W=X1, D2.W=Y1, D3.W=X2, D4.W=Y2 & Draw a line between two points \\
\ttfamily 85 & D1.W = X, D2.W = Y & Draw a line from the current pen position to (X,Y) \\
\ttfamily 86 & D1.W = X, D2.W = Y & Move the pen to (X,Y) without drawing \\
\ttfamily 87 & D1.W=X1, D2.W=Y1, D3.W=X2, D4.W=Y2 & Draw a filled rectangle \\
\ttfamily 88 & D1.W=X1, D2.W=Y1, D3.W=X2, D4.W=Y2 & Draw a filled ellipse (bounding box) \\
\ttfamily 89 & D1.W = X, D2.W = Y & Flood-fill starting at (X,Y) \\
\ttfamily 90 & D1.W=X1, D2.W=Y1, D3.W=X2, D4.W=Y2 & Draw an unfilled (outline only) rectangle \\
\ttfamily 91 & D1.W=X1, D2.W=Y1, D3.W=X2, D4.W=Y2 & Draw an unfilled (outline only) ellipse \\
\ttfamily 92 & D1.L = mode (16=single-buffer, 17=double-buffer, others=pen mode) & Set the drawing/pen mode \\
\ttfamily 93 & D1.L = width & Set the pen width \\
\ttfamily 94 & — & Repaint - explicit full flip from the back buffer to the visible screen (only meaningful in double-buffer mode) \\
\ttfamily 95 & A1=NUL-terminated string, D1.W=X, D2.W=Y & Draw a string on the graphics canvas at (X,Y) \\
\ttfamily 96 & D1.W (out) = X, D2.W (out) = Y & Return the current pen position \\
\bottomrule
\end{longtable}
\subsection*{NETWORK / SERIAL (unimplemented, by choice)}
\par\noindent\textit{Serial (40-43) and networking (100-107) log a one-time explanation and report a definite failure code, so a program checking the result can detect it and continue gracefully.}\par\medskip
\subsection*{WEB PORT EXTENSIONS (200-255 reserved)}
\par\noindent\textit{Tasks 200-255 are reserved exclusively for this web port's own debugging/convenience additions, so they can never collide with a genuine task number the real EASy68K might add.}\par\medskip
\par\noindent\textit{Tasks 200/201/203/206/213/214/215/217/220 are the CONSOLE\_LOG family: exact replicas of the real text-printing tasks 0/1/3/6/13/14/15/17/20 above (same task number + 200), except they write to the Log panel (in the same distinctive amber color) instead of the screen/output window. The real tasks 0/1/3/6/13/14/15/17/20 themselves are completely unaffected.}\par\medskip
\begin{longtable}{@{}p{1.7cm}p{5.2cm}p{7.3cm}@{}}
\toprule
\textbf{Task} & \textbf{Parameters} & \textbf{Action} \\
\midrule
\endfirsthead
\toprule
\textbf{Task} & \textbf{Parameters} & \textbf{Action} \\
\midrule
\endhead
\ttfamily 200 & A1 = string, D1.W = length & CONSOLE\_LOG - logs D1.W characters at (A1) to the Log panel (replica of task 0). \\
\ttfamily 201 & A1 = string, D1.W = length & CONSOLE\_LOG - logs D1.W characters at (A1) to the Log panel (replica of task 1). \\
\ttfamily 203 & D1.L = number & CONSOLE\_LOG - logs the decimal value of D1.L to the Log panel (replica of task 3). \\
\ttfamily 206 & D1.B = character & CONSOLE\_LOG - logs the single character in D1.B to the Log panel (replica of task 6). \\
\ttfamily 213 & A1 = NUL-terminated string & CONSOLE\_LOG - logs the string at (A1) to the Log panel (replica of task 13). \\
\ttfamily 214 & A1 = NUL-terminated string & CONSOLE\_LOG - logs the string at (A1) to the Log panel (replica of task 14). \\
\ttfamily 215 & D1.L = number, D2.B = radix (2-36) & CONSOLE\_LOG - logs D1.L converted to the given radix to the Log panel (replica of task 15). \\
\ttfamily 217 & A1 = NUL-terminated string, D1.L = value & CONSOLE\_LOG - logs the string at (A1) followed by D1.L as a decimal number, on one Log panel line (replica of task 17). \\
\ttfamily 220 & D1.L = number, D2.B = field width & CONSOLE\_LOG - logs D1.L right-justified in a D2.B-wide field to the Log panel (replica of task 20). \\
\bottomrule
\end{longtable}
\par\noindent\textit{This table covers every TRAP \#15 task as implemented in this web port. For background on TRAP \#15 itself, or any real-EASy68K behavior this table doesn't spell out, see EASy68K's own official documentation at \href{http://www.easy68k.com/}{easy68k.com}.}\par\medskip